% problem_id: S001-proposition-characterization-geometrically-regular
% source_id: S001 (Stacks Project)
% source_file: more-algebra.tex
% statement_locator: more-algebra.tex lines 8696--8711
% proof_locator: more-algebra.tex lines 8713--8848
% env: proposition
% id_note: label 在本来源内唯一
% pairing: tight (verified)
% status: complete (statement + author proof verbatim + reason + locators)
%
\begin{proposition}
\label{proposition-characterization-geometrically-regular}
Let $k$ be a field of characteristic $p > 0$.
Let $(A, \mathfrak m, K)$ be a Noetherian local
$k$-algebra. The following are equivalent
\begin{enumerate}
\item $A$ is geometrically regular over $k$,
\item for all $k \subset k' \subset k^{1/p}$
finite over $k$ the ring $A \otimes_k k'$ is regular,
\item $A$ is regular and the canonical map
$H_1(L_{K/k}) \to \mathfrak m/\mathfrak m^2$ is injective, and
\item $A$ is regular and the map
$\Omega_{k/\mathbf{F}_p} \otimes_k K \to \Omega_{A/\mathbf{F}_p} \otimes_A K$
is injective.
\end{enumerate}
\end{proposition}

\begin{proof}
Proof of (3) $\Rightarrow$ (1).
Assume (3). Let $k'/k$ be a finite purely inseparable extension.
Set $A' = A \otimes_k k'$. This is a local ring with maximal ideal
$\mathfrak m'$. Set $K' = A'/\mathfrak m'$. We get a commutative
diagram
$$
\xymatrix{
0 \ar[r] &
H_1(L_{K/k}) \otimes K' \ar[r] \ar[d]_\beta &
\mathfrak m/\mathfrak m^2 \otimes K' \ar[r] \ar[d] &
\Omega_{A/k} \otimes_A K' \ar[r] \ar[d]_{\cong} &
\Omega_{K/k} \otimes K' \ar[r] \ar[d]_\alpha &
0
\\
 &
H_1(L_{K'/k'}) \ar[r] &
\mathfrak m'/(\mathfrak m')^2 \ar[r] &
\Omega_{A'/k'} \otimes_{A'} K' \ar[r] &
\Omega_{K'/k'} \ar[r] &
0
}
$$
with exact rows. The third vertical arrow is an isomorphism by base
change for modules of differentials
(Algebra, Lemma \ref{algebra-lemma-differentials-base-change}).
Thus $\alpha$ is surjective. By
Lemma \ref{lemma-gamma-commutative-diagram} we have
$$
\dim \Ker(\alpha) - \dim \Ker(\beta) + \dim \Coker(\beta) = 0
$$
(and these dimensions are all finite). A diagram chase shows that
$\dim \mathfrak m'/(\mathfrak m')^2 \leq \dim \mathfrak m/\mathfrak m^2$.
However, since $A \to A'$ is finite flat we see that
$\dim(A) = \dim(A')$, see
Algebra, Lemma \ref{algebra-lemma-dimension-base-fibre-total}.
Hence $A'$ is regular by definition.

\medskip\noindent
Equivalence of (3) and (4). Consider the Jacobi-Zariski sequences
for rows of the commutative diagram
$$
\xymatrix{
\mathbf{F}_p \ar[r] & A \ar[r] & K \\
\mathbf{F}_p \ar[r] \ar[u] & k \ar[r] \ar[u] & K \ar[u]
}
$$
to get a commutative diagram
$$
\xymatrix{
0 \ar[r] &
\mathfrak m/\mathfrak m^2 \ar[r] &
\Omega_{A/\mathbf{F}_p} \otimes_A K \ar[r] &
\Omega_{K/\mathbf{F}_p} \ar[r] & 0 & \\
0 \ar[r] &
H_1(L_{K/k}) \ar[r] \ar[u] &
\Omega_{k/\mathbf{F}_p} \otimes_k K \ar[r] \ar[u] &
\Omega_{K/\mathbf{F}_p} \ar[r] \ar[u] &
\Omega_{K/k} \ar[r] \ar[u] &
0
}
$$
with exact rows. We have used that $H_1(L_{K/A}) = \mathfrak m/\mathfrak m^2$
and that $H_1(L_{K/\mathbf{F}_p}) = 0$ as $K/\mathbf{F}_p$ is separable, see
Algebra, Proposition
\ref{algebra-proposition-characterize-separable-field-extensions}.
Thus it is clear that the kernels of
$H_1(L_{K/k}) \to \mathfrak m/\mathfrak m^2$
and
$\Omega_{k/\mathbf{F}_p} \otimes_k K \to \Omega_{A/\mathbf{F}_p} \otimes_A K$
have the same dimension.

\medskip\noindent
Proof of (2) $\Rightarrow$ (4) following Faltings, see
\cite{Faltings-einfacher}. Let $a_1, \ldots, a_n \in k$ be elements
such that $\text{d}a_1, \ldots, \text{d}a_n$ are linearly independent
in $\Omega_{k/\mathbf{F}_p}$. Consider the field extension
$k' = k(a_1^{1/p}, \ldots, a_n^{1/p})$. By
Algebra, Lemma \ref{algebra-lemma-size-extension-pth-roots}
we see that $k' = k[x_1, \ldots, x_n]/(x_1^p - a_1, \ldots, x_n^p - a_n)$.
In particular we see that the naive cotangent complex of $k'/k$
is homotopic to the complex
$\bigoplus_{j = 1, \ldots, n} k' \rightarrow \bigoplus_{i = 1, \ldots, n} k'$
with the zero differential as
$\text{d}(x_j^p - a_j) = 0$ in $\Omega_{k[x_1, \ldots, x_n]/k}$.
Set $A' = A \otimes_k k'$ and $K' = A'/\mathfrak m'$ as above.
By Algebra, Lemma \ref{algebra-lemma-change-base-NL}
we see that $\NL_{A'/A}$ is homotopy equivalent to the complex
$\bigoplus_{j = 1, \ldots, n} A' \rightarrow \bigoplus_{i = 1, \ldots, n} A'$
with the zero differential, i.e., $H_1(L_{A'/A})$ and
$\Omega_{A'/A}$ are free of rank $n$. The Jacobi-Zariski sequence for
$\mathbf{F}_p \to A \to A'$ is
$$
H_1(L_{A'/A}) \to \Omega_{A/\mathbf{F}_p} \otimes_A A'
\to \Omega_{A'/\mathbf{F}_p} \to \Omega_{A'/A} \to 0
$$
Using the presentation $A[x_1, \ldots, x_n] \to A'$ with
kernel $(x_j^p - a_j)$ we see, unwinding the maps in
Algebra, Lemma \ref{algebra-lemma-exact-sequence-NL},
that the $j$th basis vector of $H_1(L_{A'/A})$ maps to
$\text{d}a_j \otimes 1$ in $\Omega_{A/\mathbf{F}_p} \otimes A'$.
As $\Omega_{A'/A}$ is free (hence flat) we get on tensoring with $K'$
an exact sequence
$$
K'^{\oplus n} \to \Omega_{A/\mathbf{F}_p} \otimes_A K'
\xrightarrow{\beta} \Omega_{A'/\mathbf{F}_p} \otimes_{A'} K' \to
K'^{\oplus n} \to 0
$$
We conclude that the elements $\text{d}a_j \otimes 1$ generate
$\Ker(\beta)$ and we have to show that are linearly independent, i.e.,
we have to show $\dim(\Ker(\beta)) = n$.
Consider the following big diagram
$$
\xymatrix{
0 \ar[r] &
\mathfrak m'/(\mathfrak m')^2 \ar[r] &
\Omega_{A'/\mathbf{F}_p} \otimes K' \ar[r] &
\Omega_{K'/\mathbf{F}_p} \ar[r] & 0 \\
0 \ar[r] &
\mathfrak m/\mathfrak m^2 \otimes K' \ar[r] \ar[u]^\alpha &
\Omega_{A/\mathbf{F}_p} \otimes K' \ar[r] \ar[u]^\beta &
\Omega_{K/\mathbf{F}_p} \otimes K' \ar[r] \ar[u]^\gamma & 0
}
$$
By Lemma \ref{lemma-cartier-equality} and the Jacobi-Zariski sequence for
$\mathbf{F}_p \to K \to K'$ we see that the kernel and cokernel of
$\gamma$ have the same finite dimension.
By assumption $A'$ is regular (and of the same dimension as $A$, see
above) hence the kernel and cokernel of $\alpha$ have the same dimension.
It follows that the kernel and cokernel of $\beta$ have the same
dimension which is what we wanted to show.

\medskip\noindent
The implication (1) $\Rightarrow$ (2) is trivial. This finishes
the proof of the proposition.
\end{proof}
